\begin{textsample}{POS Dim 1 – 2009 – Score 24.00 – 2009/t002945.txt}  \label{ex:f1_pos_030}
\textbf{The} substance of things \textbf{unseen}.
Cities, past and future.
In Oxford, perhaps we can use Lewis Carroll and look in \textbf{the} looking glass that is New York City to try and see our true selves, or perhaps pass through to another world.
Or, in \textbf{the} words of F.
Scott Fitzgerald,"As \textbf{the} moon rose higher, \textbf{the} inessential houses began to melt away until gradually I became aware of \textbf{the} old \textbf{island} here that once flowered for Dutch sailors ' eyes, a fresh green breast of \textbf{the} new world."My colleagues and I have been working for 10 years to rediscover this lost world in a project we call \textbf{The} Mannahatta Project.
We ' re trying to discover what Henry Hudson would have seen on \textbf{the} afternoon of September 12th, 1609, when he sailed into New York harbor.
And I ' d like to tell you \textbf{the} story in three acts, and if I have time still, an epilogue.
So, Act I : A Map Found.
So, I didn ' t grow up in New York.
I grew up out west in \textbf{the} Sierra Nevada Mountains, like you see here, in \textbf{the} Red Rock Canyon.
And from these early experiences as a child I learned to love landscapes.
And so when it became time for me to do my graduate studies, I studied this emerging field of landscape \textbf{ecology}.
Landscape \textbf{ecology} concerns itself with how \textbf{the} stream and \textbf{the} meadow and \textbf{the} forest and \textbf{the} cliffs make \textbf{habitats} for plants and animals.
This experience and this training lead me to get a wonderful job with \textbf{the} Wildlife Conservation Society, which works to save wildlife and wild places all over \textbf{the} world.
And over \textbf{the} last decade, I traveled to over 40 countries to see jaguars and bears and \textbf{elephants} and \textbf{tigers} and rhinos.
But every time I would return from my trips I ' d return back to New York City.
And on my weekends I would go up, just like all \textbf{the} other tourists, to \textbf{the} top of \textbf{the} Empire State Building, and I ' d look down on this landscape, on these ecosystems, and I ' d wonder,"How does this landscape work to make \textbf{habitat} for plants and animals?
How does it work to make \textbf{habitat} for animals like me?"I ' d go to Times Square and I ' d look at \textbf{the} amazing ladies on \textbf{the} wall, and wonder why nobody is looking at \textbf{the} historical figures just behind them.
I ' d go to Central Park and see \textbf{the} rolling topography of Central Park come up against \textbf{the} abrupt and \textbf{sheer} topography of midtown Manhattan.
I started reading about \textbf{the} history and \textbf{the} geography in New York City.
I read that New York City was \textbf{the} first mega-city, a city of 10 million people or more, in 1950.
I started seeing paintings like this.
For those of you who are from New York, this is 125th street under \textbf{the} West Side Highway.
It was once a beach.
And this painting has John James Audubon, \textbf{the} painter, sitting on \textbf{the} \textbf{rock}.
And it ' s looking up on \textbf{the} wooded heights of Washington Heights to Jeffrey ' s Hook, where \textbf{the} George Washington Bridge goes across today.
Or this painting, from \textbf{the} 1740s, from Greenwich Village.
Those are two students at King ' s College — later Columbia University — sitting on a hill, overlooking a valley.
And so I ' d go down to Greenwich Village and I ' d look for this hill, and I couldn ' t find it.
And I couldn ' t find that palm tree.
What ' s that palm tree doing there?
So, it was in \textbf{the} course of these investigations that I ran into a map.
And it ' s this map you see here.
It ' s held in a geographic information system which allows me to zoom in.
This map isn ' t from Hudson ' s time, but from \textbf{the} American Revolution, 170 years later, made by British military cartographers during \textbf{the} occupation of New York City.
And it ' s a remarkable map.
It ' s in \textbf{the} National Archives here in Kew.
And it ' s 10 feet long and three and a half feet wide.
And if I zoom in to lower Manhattan you can see \textbf{the} extent of New York City as it was, right at \textbf{the} end of \textbf{the} American Revolution.
Here ' s Bowling Green.
And here ' s Broadway.
And this is City Hall Park.
So \textbf{the} city basically extended to City Hall Park.
And just beyond it you can see features that have vanished, things that have \textbf{disappeared}.
This is \textbf{the} Collect Pond, which was \textbf{the} fresh water source for New York City for its first 200 years, and for \textbf{the} Native Americans for thousands of years before that.
You can see \textbf{the} Lispenard Meadows draining down through here, through what is TriBeCa now, and \textbf{the} beaches that come up from \textbf{the} Battery, all \textbf{the} way to 42nd St.
This map was made for military reasons.
They ' re mapping \textbf{the} roads, \textbf{the} buildings, these fortifications that they built.
But they ' re also mapping things of ecological interest, also military interest : \textbf{the} hills, \textbf{the} marshes, \textbf{the} streams.
This is Richmond Hill, and Minetta Water, which used to run its way through Greenwich Village.
Or \textbf{the} swamp at Gramercy Park, right here.
Or Murray Hill.
And this is \textbf{the} Murrays ' house on Murray Hill, 200 years ago.
Here is Times Square, \textbf{the} two streams that came together to make a wetland in Times Square, as it was at \textbf{the} end of \textbf{the} American Revolution.
So I saw this remarkable map in a book.
And I thought to myself,"You know, if I could georeference this map, if I could place this map in \textbf{the} grid of \textbf{the} city today, I could find these lost features of \textbf{the} city, in \textbf{the} block-by-block geography that people know, \textbf{the} geography of where people go to work, and where they go to live, and where they like to eat."So, after some work we were able to georeference it, which allows us to put \textbf{the} modern streets on \textbf{the} city, and \textbf{the} buildings, and \textbf{the} open spaces, so that we can zoom in to where \textbf{the} Collect Pond is.
We can digitize \textbf{the} Collect Pond and \textbf{the} streams, and see where they actually are in \textbf{the} geography of \textbf{the} city today.
So this is fun for finding where things are \textbf{relative} to \textbf{the} old topography.
But I had another idea about this map.
If we take away \textbf{the} streets, and if we take away \textbf{the} buildings, and if we take away \textbf{the} open spaces, then we could take this map.
If we pull off \textbf{the} 18th century features we could drive it back in time.
We could drive it back to its ecological fundamentals : to \textbf{the} hills, to \textbf{the} streams, to \textbf{the} basic hydrology and shoreline, to \textbf{the} beaches, \textbf{the} basic aspects that make \textbf{the} ecological landscape.
Then, if we added maps like \textbf{the} geology, \textbf{the} bedrock geology, and \textbf{the} \textbf{surface} geology, what \textbf{the} glaciers leave, if we make \textbf{the} soil map, with \textbf{the} 17 soil classes, that are defined by \textbf{the} National Conservation Service, if we make a digital elevation model of \textbf{the} topography that tells us how high \textbf{the} hills were, then we can \textbf{calculate} \textbf{the} \textbf{slopes}.
We can \textbf{calculate} \textbf{the} aspect.
We can \textbf{calculate} \textbf{the} winter wind exposure — so, which way \textbf{the} winter winds blow across \textbf{the} landscape. \textbf{The} white areas on this map are \textbf{the} places protected from \textbf{the} winter winds.
We compiled all \textbf{the} information about where \textbf{the} Native Americans were, \textbf{the} Lenape.
And we built a \textbf{probability} map of where they might have been.
So, \textbf{the} red areas on this map indicate \textbf{the} places that are best for human \textbf{sustainability} on Manhattan, places that are close to water, places that are near \textbf{the} harbor to \textbf{fish}, places protected from \textbf{the} winter winds.
We know that there was a Lenape settlement down here by \textbf{the} Collect Pond.
And we knew that they planted a kind of horticulture, that they grew these beautiful gardens of corn, beans, and squash, \textbf{the}"Three Sisters"garden.
So, we built a model that explains where those fields might have been.
And \textbf{the} old fields, \textbf{the} successional fields that go.
And we might think of these as abandoned.
But, in fact, they ' re grassland \textbf{habitats} for grassland birds and plants.
And they have become successional shrub lands, and these then mix in to a map of all \textbf{the} ecological communities.
And it turns out that Manhattan had 55 different ecosystem types.
You can think of these as neighborhoods, as distinctive as TriBeCa and \textbf{the} Upper East Side and Inwood — that these are \textbf{the} forest and \textbf{the} wetlands and \textbf{the} marine communities, \textbf{the} beaches.
And 55 is a lot.
On a per-area basis, Manhattan had more ecological communities per acre than Yosemite does, than Yellowstone, than Amboseli.
It was really an extraordinary landscape that was capable of supporting an extraordinary biodiversity.
So, Act II : A Home Reconstructed.
So, we studied \textbf{the} \textbf{fish} and \textbf{the} frogs and \textbf{the} birds and \textbf{the} bees, \textbf{the} 85 different kinds of \textbf{fish} that were on Manhattan, \textbf{the} Heath hens, \textbf{the} species that aren ' t there anymore, \textbf{the} beavers on all \textbf{the} streams, \textbf{the} black bears, and \textbf{the} Native Americans, to study how they used and thought about their landscape.
We wanted to try and map these.
And to do that what we did was we mapped their \textbf{habitat} needs.
Where do they get their food?
Where do they get their water?
Where do they get their shelter?
Where do they get their reproductive resources?
To an ecologist, \textbf{the} intersection of these is \textbf{habitat}, but to most people, \textbf{the} intersection of these is their home.
So, we would read in field guides, \textbf{the} standard field guides that maybe you have on your shelves, you know, what beavers need is,"A slowly meandering stream with aspen trees and alders and willows, near \textbf{the} water."That ' s \textbf{the} best thing for a beaver.
So we just started making a list.
Here is \textbf{the} beaver.
And here is \textbf{the} stream, and \textbf{the} aspen and \textbf{the} alder and \textbf{the} willow.
As if these were \textbf{the} maps that we would need to predict where you would find \textbf{the} beaver.
Or \textbf{the} bog \textbf{turtle}, needing \textbf{wet} meadows and insects and sunny places.
Or \textbf{the} bobcat, needing rabbits and beavers and den sites.
And rapidly we started to realize that beavers can be something that a bobcat needs.
But a beaver also needs things.
And that having it on either side means that we can link it together, that we can create \textbf{the} network of \textbf{the} \textbf{habitat} relationships for these species.
Moreover, we realized that you can start out as being a beaver specialist, but you can look up what an aspen needs.
An aspen needs fire and dry soils.
And you can look at what a \textbf{wet} meadow needs.
And it need beavers to create \textbf{the} wetlands, and maybe some other things.
But you can also talk about sunny places.
So, what does a sunny place need?
Not \textbf{habitat} per se.
But what are \textbf{the} conditions that make it possible?
Or fire.
Or dry soils.
And that you can put these on a grid that ' s 1, 000 columns long across \textbf{the} top and 1, 000 rows down \textbf{the} other way.
And then we can visualize this data like a network, like a social network.
And this is \textbf{the} network of all \textbf{the} \textbf{habitat} relationships of all \textbf{the} plants and animals on Manhattan, and everything they needed, going back to \textbf{the} geology, going back to time and space at \textbf{the} very core of \textbf{the} \textbf{web}.
We call this \textbf{the} Muir Web.
And if you zoom in on it it looks like this.
Each point is a different species or a different stream or a different soil type.
And those little gray lines are \textbf{the} connections that connect them together.
They are \textbf{the} connections that actually make nature resilient.
And \textbf{the} structure of this is what makes nature work, seen with all its parts.
We call these Muir \textbf{Webs} after \textbf{the} Scottish-American naturalist John Muir, who said,"When we try to pick out anything by itself, we find that it ' s bound fast by a thousand invisible cords that cannot be broken, to everything in \textbf{the} \textbf{universe}."So then we took \textbf{the} Muir \textbf{webs} and we took them back to \textbf{the} maps.
So if we wanted to go between 85th and 86th, and Lex and Third, maybe there was a stream in that block.
And these would be \textbf{the} kind of trees that might have been there, and \textbf{the} flowers and \textbf{the} lichens and \textbf{the} mosses, \textbf{the} butterflies, \textbf{the} \textbf{fish} in \textbf{the} stream, \textbf{the} birds in \textbf{the} trees.
Maybe a timber rattlesnake lived there.
And perhaps a black bear walked by.
And maybe Native Americans were there.
And then we took this data.
You can see this for yourself on our website.
You can zoom into any block on Manhattan, and see what might have been there 400 years ago.
And we used it to try and reveal a landscape here in Act III.
We used \textbf{the} tools they use in Hollywood to make these \textbf{fantastic} landscapes that we all see in \textbf{the} movies.
And we tried to use it to visualize Third Avenue.
So we would take \textbf{the} landscape and we would build up \textbf{the} topography.
We ' d lay on top of that \textbf{the} soils and \textbf{the} waters, and illuminate \textbf{the} landscape.
We would lay on top of that \textbf{the} map of \textbf{the} ecological communities.
And feed into that \textbf{the} map of \textbf{the} species.
So that we would actually take a photograph, flying above Times Square, looking toward \textbf{the} Hudson River, waiting for Hudson to come.
Using this technology, we can make these \textbf{fantastic} georeferenced views.
We can basically take a picture out of any window on Manhattan and see what that landscape looked like 400 years ago.
This is \textbf{the} view from \textbf{the} East River, looking up Murray Hill at where \textbf{the} United Nations is today.
This is \textbf{the} view looking down \textbf{the} Hudson River, with Manhattan on \textbf{the} \textbf{left}, and New Jersey out on \textbf{the} right, looking out toward \textbf{the} Atlantic Ocean.
This is \textbf{the} view over Times Square, with \textbf{the} beaver pond there, looking out toward \textbf{the} east.
So we can see \textbf{the} Collect Pond, and Lispenard Marshes back behind.
We can see \textbf{the} fields that \textbf{the} Native Americans made.
And we can see this in \textbf{the} geography of \textbf{the} city today.
So when you ' re watching"Law and Order,"and \textbf{the} lawyers walk up \textbf{the} steps they could have walked back down those steps of \textbf{the} New York Court House, right into \textbf{the} Collect Pond, 400 years ago.
So these images are \textbf{the} work of my friend and colleague, Mark Boyer, who is here in \textbf{the} audience today.
And I ' d just like, if you would give him a hand, to call out for his fine work.
There is such power in bringing science and visualization together, that we can create images like this, perhaps looking on either side of a looking glass.
And even though I ' ve only had a brief time to speak, I hope you appreciate that Mannahatta was a very special place. \textbf{The} place that you see here on \textbf{the} \textbf{left} side was interconnected.
It was based on this diversity.
It had this resilience that is what we need in our modern world.
But I wouldn ' t have you think that I don ' t like \textbf{the} place on \textbf{the} right, which I quite do.
I ' ve come to love \textbf{the} city and its kind of diversity, and its resilience, and its dependence on density and how we ' re connected together.
In fact, that I see them as reflections of each other, much as Lewis Carroll did in"Through \textbf{the} Looking Glass."We can compare these two and hold them in our minds at \textbf{the} same time, that they really are \textbf{the} same place, that there is no way that cities can escape from nature.
And I think this is what we ' re learning about building cities in \textbf{the} future.
So if you ' ll allow me a brief epilogue, not about \textbf{the} past, but about 400 years from now, what we ' re realizing is that cities are \textbf{habitats} for people, and need to supply what people need : a sense of home, food, water, shelter, reproductive resources, and a sense of meaning.
This is \textbf{the} particular additional \textbf{habitat} requirement of humanity.
And so many of \textbf{the} talks here at TED are about meaning, about bringing meaning to our lives in all kinds of different ways, through technology, through art, through science, so much so that I think we focus so much on that side of our lives, that we haven ' t given enough attention to \textbf{the} food and \textbf{the} water and \textbf{the} shelter, and what we need to raise \textbf{the} kids.
So, how can we envision \textbf{the} city of \textbf{the} future?
Well, what if we go to Madison Square Park, and we imagine it without all \textbf{the} cars, and bicycles instead and large forests, and streams instead of sewers and storm drains?
What if we imagined \textbf{the} Upper East Side with green roofs, and streams winding through \textbf{the} city, and windmills supplying \textbf{the} power we need?
Or if we imagine \textbf{the} New York City metropolitan area, currently home to 12 million people, but 12 million people in \textbf{the} future, perhaps living at \textbf{the} density of Manhattan, in only 36 percent of \textbf{the} area, with \textbf{the} areas in between covered by farmland, covered by wetlands, covered by \textbf{the} marshes we need.
This is \textbf{the} kind of future I think we need, is a future that has \textbf{the} same diversity and abundance and dynamism of Manhattan, but that learns from \textbf{the} \textbf{sustainability} of \textbf{the} past, of \textbf{the} \textbf{ecology}, \textbf{the} original \textbf{ecology}, of nature with all its parts.
Thank you very much.

% matched lemmas: calculate, disappear, ecology, elephant, fantastic, fish, habitat, island, left, probability, relative, rock, sheer, slope, surface, sustainability, the, tiger, turtle, universe, unseen, web, wet
\end{textsample}
